\begingroup
\def\C{\mathbb{C}}
\def\R{\mathbb{R}}
\def\Tr{\operatorname{Tr}}
\def\diag{\operatorname{diag}}
\section{Exact endpoint windows: Gamma coefficients, paired Schur volumes, and intrinsic sensitivity}



This continuation evaluates the original endpoint quantity of the complete
written criterion at source revision
\nolinkurl{c720f40530eed2f5969dbabe94dfd3fddc0f507f} of PR~23.
The source criterion and the supplied \emph{Endpoint degree selection in
the actual gamma-corrected coordinates} were read completely, together
with GC.1--59 and GMT.1--42, including GMT.29a--b. The present work proves
the arbitrary-window Schur map, its exact determinant cancellations, a
finite rational interval certificate, and a new sensitivity estimate
through the actual angles between canonical least-lift spaces. The
finite input list grows through degree $2(n+r)$. Every statement below
is finite. No estimate for that growing list or its arithmetic
conditioning is presumed.

\subsection{The original source and the complete finite input map}

Let $h(s)$ be the given monic packet of actual zeros, with their complete
orders, of
\[
 g(s)=s(s-1)\pi^{-s/2}\Gamma(s/2)\zeta(s)=2\xi(s).
\]
Retain its entire quotient $v_h=g/h$ and original positive density
\[
 w_h(t)=\frac{|v_h(1/2+it)|^2}{2\pi},\qquad
 \mu_h=\int_{\R} w_h(t)\,dt,\qquad m_{h,k}=w_h^{*k}.
                                                               \tag{EW.1}
\]
The complete analytic construction in GC.2--9 proves that $w_h$ has
every polynomial moment and exponential moments of every order less
than $\pi/2$. Its mechanism is the theta seed
$\phi_0=(4\pi^2x^4-6\pi x^2)e^{-\pi x^2}$, the equality
$\mathcal M\Theta\phi_0=g$, contour rotation through every angle
$|\beta|<\pi/4$, and exact local division by $h$.
In particular $|v_h(\sigma+it)|$ is bounded, uniformly for $\sigma$
in each fixed compact real interval, by
$C_{j,\beta,h}(1+|t|)^{-j}e^{-\beta|t|}$ for every
$j\ge0$ and $0<\beta<\pi/4$. These proofs and their original
constants are retained in the accompanying complete GC source.
The zero set of $v_h(1/2+it)$ is discrete, so $w_h>0$ almost
everywhere. Its convolution has positive mass $\mu_h^k$.

Throughout this paper
\[
 \begin{gathered}
 k\ge1,\quad c=k/2,\quad S=s_1+\cdots+s_k=c+iu,
 \quad u=t_1+\cdots+t_k,\\
 \langle P,Q\rangle_H=
 \int_{\R}\overline{P(c+iu)}Q(c+iu)m_{h,k}(u)\,du.
 \end{gathered}                                               \tag{EW.2}
\]
The pairing is conjugate-linear in its first argument. The sum chart
$(t_1,\ldots,t_k)\mapsto(t_1,\ldots,t_{k-1},u)$ has determinant
one. Fubini therefore proves both the convolution in (EW.1) and the
full mass $\mu_h^k$, with no average coordinate or mass division.
Every nonzero polynomial has finitely many real-line zeros, so its
squared norm in (EW.2) is strictly positive.

Fix the original monic relation $\chi\in\C[S]$, of degree $q\ge1$.
For the packet source it generates the exact intersection
$(h(s_1),\ldots,h(s_k))\cap\C[S]$ under $S=\sum s_i$.
Define $\mathcal C=\C[S]/(\chi)$ in the ordered coordinates
$1,S,\ldots,S^{q-1}$. The quotient map is the actual remainder
map. Its injection into the tensor arithmetic jets is
\[
 \eta:\mathcal C\longrightarrow(\C[s]/(h))^{\otimes k},
 \qquad [P]\longmapsto[v_h]_h^{\otimes k}P(s_1+\cdots+s_k).
                                                               \tag{EW.3}
\]
Complete local cancellation makes the constant term of $v_h$ nonzero
at every selected centre. Each local truncated power series is
invertible by recursively solving its product with an inverse to
equal one. The Chinese remainder isomorphism then proves that
$[v_h]_h$ is a unit. The unfactored polynomial map
$\C[S]\to(\C[s]/(h))^{\otimes k}$ in (EW.3) consequently has
kernel equal to the stated intersection ideal $(\chi)$. Factoring
by that ideal gives the displayed $\eta$ with kernel zero,
proving injectivity without discarding a jet.

For $N\ge q-1$, put $\mathcal P_N=\C[S]_{\le N}$ and
$\mathcal P_{-1}=0$. Let $H_N$ be its source matrix in the original
monomials. Let $J_N:\C^{N+1}\to\C^q$ be monic remainder
and $B_N:\C^{N-q+1}\to\C^{N+1}$ be multiplication by $\chi$.
Monic division proves the exact sequence
\[
 0\longrightarrow\mathcal P_{N-q}
 \xrightarrow{M_\chi}\mathcal P_N
 \xrightarrow{J_N}\mathcal C\longrightarrow0.                \tag{EW.4}
\]
In particular $B_N$ is injective, $J_N$ is surjective, and
$\ker J_N=\operatorname{im}B_N$.

For later coefficient evaluation keep all Gamma constants:
\[
 \lambda>0,\quad\alpha=2\lambda,\quad\beta=k\alpha,\quad
 r_\lambda(t)=\frac{|\Gamma(\lambda+it/2)|^2}{2\pi},\quad
 c_\lambda=2^{1-2\lambda}\Gamma(2\lambda),\quad C_k=c_\lambda^k.
                                                               \tag{EW.5}
\]
Let the real monic polynomials $b_j^{(\lambda)}$ be specified by
\[
 \sum_{j\ge0}\frac{b_j^{(\lambda)}(t)}{j!}z^j
 =(1+z^2)^{-\lambda}e^{t\arctan z},\quad |z|<1,
 \qquad
 b_{j+1}=tb_j-j(j+\alpha-1)b_{j-1},                          \tag{EW.6}
\]
where the logarithms defining the generating function are the branches
from zero, $b_0=1$, and $b_1=t$. Differentiation of the displayed
generating function proves the recurrence and monicity. The beta
integral and Fourier inversion give the literal convolution mass and
Laplace transform
\[
 r_{\lambda,k}=r_\lambda^{*k}
   =\frac{c_\lambda^k}{c_{k\lambda}}r_{k\lambda},\qquad
 \int e^{\theta u}r_{\lambda,k}(u)\,du
    =C_k(\cos\theta)^{-\beta},\quad |\Re\theta|<\pi/2.
                                                               \tag{EW.7}
\]
Their complete beta-integral proof with the changes $t=2v$ and
$u=\sum t_i$ is GC.10--15. In particular the reference monic norms
at $\theta=0$ are $C_kj!(\beta)_j$. They also follow directly by
integrating the product of two functions (EW.6), which gives
$C_k(1-zw)^{-\beta}$, and comparing its finite derivatives.

Define the given one-factor coefficients and their finite powers by
\[
 c_{h,j}=\frac{\int b_j^{(\lambda)}(t)w_h(t)\,dt}
                  {c_\lambda j!(\alpha)_j},\qquad
 d_l=\sum_{j_1+\cdots+j_k=l}\prod_{a=1}^k(\alpha)_{j_a}c_{h,j_a}.
                                                               \tag{EW.8}
\]
Multiplying the $k$ generating functions (EW.6) and comparing the
coefficient of $z^l$ gives the finite polynomial identity
\[
 b_l^{(k\lambda)}(\textstyle\sum t_i)
   =\sum_{j_1+\cdots+j_k=l}\frac{l!}{\prod j_i!}
                      \prod_i b_{j_i}^{(\lambda)}(t_i).
\]
Integrating every term against the full product measure proves
\[
 \int b_l^{(k\lambda)}(u)m_{h,k}(u)\,du=C_kl!d_l.
                                                               \tag{EW.9}
\]
All sums here are finite and all their terms are integrable. For any
specified polynomials $p_a\in\mathcal P_M$, perform monic
polynomial subtraction in the actual $u$ variable to find the unique
coefficients $L_{ab,l}$ in
\[
 \overline{p_a(c+iu)}p_b(c+iu)
      =\sum_{l=0}^{2M}L_{ab,l}b_l^{(k\lambda)}(u).
\]
The complete coefficient-to-Gram map is therefore
\[
 H(p)_{ab}=C_k\sum_{l=0}^{2M}l!L_{ab,l}d_l.                  \tag{EW.10}
\]
It is a degree-$k$ real polynomial map from the actual coefficient
vector to Hermitian forms. The precise tests contain $c=k/2$, the
two phases $\pm i$, and, when a relation column is used, both
$\chi(c+iu)$ and $\overline\chi(c-iu)$. In particular (EW.10)
constructs $H_M$ from $c_{h,0},\ldots,c_{h,2M}$ and the stated
$k,\lambda$. At $l=0$, $c_{h,0}=\mu_h/c_\lambda$, so the
same map gives $H_{00}=\mu_h^k$.

\subsection{Canonical lifts and one arbitrary-window block}

For each admitted degree define the exact matrices and maps
\[
 K_N=J_NH_N^{-1}J_N^*,\quad G_N=K_N^{-1},\quad
 R_N=H_N^{-1}J_N^*G_N:\C^q\longrightarrow\C^{N+1},\quad
 V_N=\det G_N.                                               \tag{EW.11}
\]
Positivity of $H_N$ and surjectivity of $J_N$ prove $K_N>0$.
Multiplication gives $J_NR_N=I$ and $B_N^*H_NR_N=0$.
For every $x\in\C^q$, any lift is $R_Nx+B_Ny$, and its norm
is exactly $x^*G_Nx+y^*B_N^*H_NB_Ny$. This proves the unique
minimum property and $R_N^*H_NR_N=G_N$.

Let $E_N$ include the first $q$ monomials. The square coefficient
matrix $[E_N,B_N]$ has determinant one: its columns are ordered by
successive degree and each leading coefficient is one. Subtracting
boundary columns to replace $E_N$ by $R_N$ retains determinant
one. Its Gram is block diagonal, so
\[
 V_N=\frac{\det H_N}{\det(B_N^*H_NB_N)}.                    \tag{EW.12}
\]
Empty determinants are one. This fixes the relation determinant and
its index $N-q+1$ in the original quotient coordinates.

Take integers
\[
 n\ge q,\qquad r\ge1,\qquad m=n+r,\qquad d=r+1.
                                                               \tag{EW.13}
\]
In $\mathcal P_m$ split the ordered monomials into degrees
$0,\ldots,n-1$ and the $d$ appended columns $S^n,\ldots,S^m$.
Write
\[
 H_m=\begin{pmatrix}H_{n-1}&C_D\\C_D^*&D_D\end{pmatrix},\quad
 W=D_D-C_D^*H_{n-1}^{-1}C_D,\quad
 J_m=[J_{n-1},J_D],\quad F=J_D-J_{n-1}H_{n-1}^{-1}C_D.
                                                               \tag{EW.14}
\]
The map
\[
 \Psi:\C^n\oplus\C^d\longrightarrow\C^{m+1},\qquad
 (x,y)\longmapsto(x-H_{n-1}^{-1}C_Dy,y)                     \tag{EW.15}
\]
is invertible with determinant one. Direct multiplication gives
$\Psi^*H_m\Psi=\diag(H_{n-1},W)$ and
$J_m\Psi=[J_{n-1},F]$. Thus $W>0$; $F$ maps the appended
orthogonal residual polynomials to their original quotient classes.
No quotient is independently chosen.

Let $W_j$ be the first $j$ rows and columns of $W$, $F_j$ its
first $j$ quotient columns, and $Z_j=F_j^*G_{n-1}F_j$ for
$1\le j\le d$. The same map (EW.15), restricted to these columns,
gives the exact positive-kernel update
\[
 K_{n+j-1}=K_{n-1}+F_jW_j^{-1}F_j^*.                         \tag{EW.16}
\]
Eliminating the two diagonal blocks of
$\left(\begin{smallmatrix}I&A\\-B&I\end{smallmatrix}\right)$
in the two orders proves $\det(I+AB)=\det(I+BA)$.
Use this identity in (EW.16) to obtain
\[
 \boxed{\mathscr R_j:=\frac{V_{n-1}}{V_{n+j-1}}
     =\frac{\det(W_j+Z_j)}{\det W_j},\qquad \mathscr R_0=1.}
                                                               \tag{EW.17}
\]
All the determinants on the right are of size $j$, whereas its
equivalent kernel calculation has size $q$. The matrices $Z_j$
are positive semidefinite. Diagonalizing
$W_j^{-1/2}Z_jW_j^{-1/2}$ proves $\mathscr R_j\ge1$.
Nested minimization also proves $G_{N+1}\preceq G_N$, so
$\mathscr R_j$ is nondecreasing in $j$.

Let $p_N(S)$ be the unique monic polynomial of degree $N$
orthogonal to $\mathcal P_{N-1}$, and let $\omega_N$ be its
original squared norm. Existence and uniqueness follow by solving the
positive lower-degree projection problem; its leading coefficient
remains one. The first diagonal entry $w=W_1$ is $\omega_n$.
Eliminating the first $r$ residual columns from the final one gives
$\omega_m=\det W_d/\det W_r$. Consequently
\[
 \boxed{U:=\frac{\omega_m}{\omega_n}
    =\frac{\det W_d}{\det W_r\,w},\qquad
 \mathscr R:=\frac{V_{n-1}V_n}{V_{m-1}V_m}
    =\frac{\mathscr R_r\mathscr R_d}{\mathscr R_1}\ge1.}
                                                               \tag{EW.18}
\]
The last inequality also follows by multiplying
$V_{n-1}/V_{m-1}\ge1$ and $V_n/V_m\ge1$. Both numerator
and denominator norms carry their literal masses.

For the Gamma source (EW.7), the same $J_m$ and $\chi$ give
matrices $W^\Gamma,F^\Gamma,Z_j^\Gamma$ by the identical maps.
For their original source determinants put
$\mathfrak D_j=\det H_{j-1}$ for $j\ge1$ and
$\mathfrak D_0=1$. For the relation determinant on exactly $j$
columns put
$\mathfrak B_j=\det(B_{q+j-1}^*H_{q+j-1}B_{q+j-1})$ for
$j\ge1$ and $\mathfrak B_0=1$. The reference determinants use
the same columns and carry the superscript $\Gamma$.
Writing the supplied corrections as
$X_j=\mathfrak D_j/\mathfrak D_j^\Gamma$,
$Y_j=\mathfrak B_j/\mathfrak B_j^\Gamma$ and
$T_N=X_{N+1}/Y_{N-q+1}$, the relation (EW.12) proves
\[
 \begin{split}
 U&=(n+1)_r(n+\beta)_r
       \frac{X_{m+1}X_n}{X_mX_{n+1}},\\
 \mathscr R&=\frac{V_{n-1}^\Gamma V_n^\Gamma}
                    {V_{m-1}^\Gamma V_m^\Gamma}
             \frac{T_{n-1}T_n}{T_{m-1}T_m}.
 \end{split}                                                   \tag{EW.19}
\]
Indeed $\omega_j=C_kj!(\beta)_jX_{j+1}/X_j$ follows by
taking successive source determinants in the original monic basis.
Taking the ratio gives the first line with its exact rising factorials.
The new symbol $\mathscr R_j$ in (EW.17) has no effect on the
original positive arithmetic correction $T_N$.

\subsection{Endpoint roots and exact paired-volume cancellations}

Put $a=1/(2r)$. The original endpoint expression is
\[
 \mathcal E_{n,r}=U^a\sinh(a\log\mathscr R)
   =\frac{(U\mathscr R)^a-(U/\mathscr R)^a}{2}.              \tag{EW.20}
\]
All roots are the unique positive real roots. The equality follows
from the definition of the real hyperbolic sine and positivity.
It evaluates the existing quantity, preserving its two original norm
endpoints and all four original volume endpoints.

There is a further exact cancellation inside its two root arguments.
Set
\[
 p=\det W_r,\quad s=\det W_d,\quad
 e=w+Z_1,\quad
 D=\det(W_r+Z_r)\det(W_d+Z_d).
\]
The scalar $Z_1$ is real nonnegative, and all four quantities
$p,s,e,D$ are strictly positive. Substituting (EW.17)--(EW.18)
without changing a coordinate gives
\[
 U=\frac{s}{pw},\qquad\mathscr R=\frac{Dw}{pse},\qquad
 \boxed{U\mathscr R=\frac{D}{p^2e},\qquad
        \frac U{\mathscr R}=\frac{s^2e}{w^2D}.}              \tag{EW.21}
\]
Thus the common highest source determinant cancels from the first
argument; the common prefix source determinant cancels from the
second. The first and last original monic norms have not been set
equal or removed from the underlying source. Formula (EW.21) is the
exact rational map explaining these cancellations.

It is sometimes useful to retain the small difference in (EW.20)
as one positive numerator. With
$P=(U\mathscr R)^a$ and $Q=(U/\mathscr R)^a$, the finite
factorization $P^{2r}-Q^{2r}=(P-Q)\sum_{j=0}^{2r-1}P^{2r-1-j}Q^j$
proves
\[
 \boxed{\mathcal E_{n,r}=
 \frac{U(\mathscr R^2-1)}
 {2\mathscr R\displaystyle\sum_{j=0}^{2r-1}P^{2r-1-j}Q^j}.}
                                                               \tag{EW.22}
\]
The denominator remains positive at $\mathscr R=1$, and the
numerator is then exactly zero. This representation preserves the
sign of the endpoint quantity and avoids subtracting two enclosures
before their common positive factor has been identified.

Every entry in (EW.14) is determined by (EW.10) through degree
$2m=2(n+r)$. Each $J_N$ and $B_N$ is computed by monic division
or multiplication by the same $\chi$. Equations (EW.11)--(EW.22)
therefore prove the complete finite map
\[
 \boxed{(c_{h,0},\ldots,c_{h,2(n+r)};k,\lambda,\chi)
    \longmapsto(H_m,J_m,W,F,G_{n-1})
    \longmapsto\mathcal E_{n,r}.}                            \tag{EW.23}
\]
It consists of stated polynomial operations, inverses of matrices
whose positivity has been proved, determinants, and positive roots.

For reference the endpoint selection follows from the original
degreewise estimate by a finite argument. Write
$x_j=\tfrac12\log(V_{n+j-1}/V_{n+j+1})\ge0$ and
$A_j=\sqrt{\omega_{n+j+1}/\omega_{n+j}}$.
The local estimate is $0\le\epsilon_{n+j}\le A_j\sinh x_j$.
If an $x_j$ is zero, the minimum allowance is zero. Otherwise
$(\log\sinh x)''=-1/\sinh^2x<0$ proves finite Jensen, hence
$\prod_j\sinh x_j\le\sinh(\sum_jx_j/r)^r$.
The actual products give
$\prod_jA_j=\sqrt U$ and $2\sum_jx_j=\log\mathscr R$.
Multiplying the local estimates and taking the positive $r$th root
therefore proves
\[
 \min_{0\le j<r}\epsilon_{n+j}\le\mathcal E_{n,r}.           \tag{EW.24}
\]
The local arithmetic trace inequality and its retained loss terms are
proved in the complete source TVB.32--34a; (EW.24) is their exact
finite selection, not an added local estimate or an asymptotic claim.

\subsection{Exact projection between the canonical representative spaces}

The following result supplies more information than the separate
determinant values. Use the common ambient space
$(\mathcal P_m,\langle\cdot,\cdot\rangle_H)$, with its original
source inner product. Let $\iota_j:\C^{j+1}\to\C^{m+1}$ be
monomial inclusion. The least-lift space and its exact isometry are
\[
 L_j=\iota_jR_j\C^q\subset\mathcal P_m,\qquad
 \mathscr I_j:(\C^q,G_j)\longrightarrow L_j,\quad x\longmapsto\iota_jR_jx.
                                                               \tag{EW.25}
\]
Here $\mathscr I_j^*H_m\mathscr I_j=G_j$ proves the word isometry;
its inverse on $L_j$ is the original remainder map. The original
arithmetic correction $T_N$ retains its meaning in (EW.19).

\begin{theorem}[Canonical projection and all its moving directions]
For $q-1\le a<b\le m$, let
$\Pi_b:\mathcal P_m\to\mathcal P_m$ be the ambient orthogonal
projection with range $L_b$, including that range by its fixed
polynomial inclusion. It obeys
\[
 \boxed{\Pi_b\iota_aR_ax=\iota_bR_bx\quad(x\in\C^q).}       \tag{EW.26}
\]
Its restriction $L_a\to L_b$ is invertible. Its squared singular
values, with precisely the source inner products in (EW.25), are
the eigenvalues $\sigma_i^2$ of
\[
 G_a^{-1/2}G_bG_a^{-1/2},\qquad0<\sigma_i\le1.
                                                               \tag{EW.27}
\]
At most $s_{a,b}=\min(q,b-a)$ of these values are less than one,
and
\[
 \begin{split}
 \prod_i\sigma_i^2&=V_b/V_a,\\
 \|\Pi_a-\Pi_b\|_1
 &=2\sum_i\sqrt{1-\sigma_i^2},\\
 \|\Pi_a-\Pi_b\|_1
 &\le\min\{2\sqrt{s_{a,b}\log(V_a/V_b)},\,2s_{a,b}\}.
 \end{split}                                                   \tag{EW.28}
\]
The trace norm is for operators in the stated ambient inner product.
\end{theorem}
\begin{proof}
Both representatives in (EW.26) have remainder $x$. Their difference
belongs to $\chi\mathcal P_{b-q}$ and is therefore orthogonal to
$L_b$ by (EW.11). The second representative belongs to $L_b$;
this proves the orthogonal projection identity. The same identity
and injectivity of each lift show that the restriction is invertible.
In the quotient coordinates it is the identity map from metric
$G_a$ to metric $G_b$. The explicit maps from Euclidean coordinates
are $u\mapsto\iota_aR_aG_a^{-1/2}u$ and
$v\mapsto\iota_bR_bG_b^{-1/2}v$. Their matrices give the squared
singular values (EW.27), and no original form or mass has been
changed. A projection is a contraction, so these values lie in
$(0,1]$. Taking its determinant proves the first line of (EW.28).

The boundary space at degree $b$ extends the one at degree $a$ by
exactly $b-a$ independent monic columns. Vectors in $L_a$ are
already orthogonal to the older boundary. Imposing orthogonality to
these extra $b-a$ columns has rank at most $b-a$. Its kernel is
$L_a\cap L_b$: any vector in that kernel is in $\mathcal P_b$
and orthogonal to its full boundary, hence is a minimum representative.
The intersection therefore has dimension at least $q-(b-a)$ when
this number is positive. A unit vector in $L_a$ is unchanged in
norm by $\Pi_b$ precisely when it belongs to $L_b$, by orthogonal
Pythagoras. Thus at most $s_{a,b}$ singular values are less than one.

For a complete trace-norm calculation choose orthonormal eigenvectors
$x_i\in L_a$ for the squared restriction and put
$y_i=\Pi_bx_i/\sigma_i\in L_b$. Orthogonality gives
$\langle x_i,y_j\rangle=\sigma_j\delta_{ij}$ and the $y_i$
form an orthonormal basis of $L_b$. If $\sigma_i=1$, then
$y_i=x_i$. If $\sigma_i<1$, set
$z_i=(y_i-\sigma_ix_i)/\sqrt{1-\sigma_i^2}$.
Directly taking all these inner products proves that the planes
$\operatorname{span}(x_i,z_i)$ are mutually orthogonal and
orthogonal to the common directions. In such a plane the matrix of
$\Pi_a-\Pi_b$ is exactly
\[
 \begin{pmatrix}
 1-\sigma_i^2&-\sigma_i\sqrt{1-\sigma_i^2}\\
 -\sigma_i\sqrt{1-\sigma_i^2}&-(1-\sigma_i^2)
 \end{pmatrix}.
\]
Its eigenvalues are $\pm\sqrt{1-\sigma_i^2}$. On the common
directions and the orthogonal complement of $L_a+L_b$ the
difference is zero. This proves the exact trace-norm sum.
Finally $1-x\le-\log x$ for $0<x\le1$, because the
derivative of $x-1-\log x$ is $1-1/x$ and its value at one is
zero. Cauchy--Schwarz on the at most $s_{a,b}$ nonzero terms gives
the first bound in (EW.28); bounding each square root by one gives
the second bound independently.
\end{proof}

\subsection{Window sensitivity and exact cancellation of mass changes}

The next calculation holds for every continuously differentiable family of positive
source matrices $H_m(t)$, using all their principal submatrices and
the fixed $\chi,J_j$. Thus it applies both to the actual coefficient
forms and to a segment between their certified finite enclosures.
Write $E=\dot H_m$ and define the real spectral spread
\[
 \mathfrak o_H(E)=\lambda_{\max}(H_m^{-1/2}EH_m^{-1/2})
                -\lambda_{\min}(H_m^{-1/2}EH_m^{-1/2})\ge0.
                                                               \tag{EW.29}
\]
The two congruences in this expression specify its exact coordinate
map; the source $H_m$ remains in every evaluation.

Embed the coefficient vector of $p_j$ in $\mathcal P_m$ and set
\[
 P_j=\frac{\iota_jp_jp_j^*\iota_j^*}{\omega_j},\qquad
 Q_j=\iota_jR_jG_j^{-1}R_j^*\iota_j^*.
                                                               \tag{EW.30}
\]
These are contravariant matrices; their products $P_jH_m$ and
$Q_jH_m$ are respectively the ambient orthogonal projections onto
the monic orthogonal line and $L_j$. To verify this, multiply each
by itself and use $p_j^*H_jp_j=\omega_j$ or
$R_j^*H_jR_j=G_j$. The resulting operators are idempotent and
self-adjoint in $H_m$ and have the stated ranges. The matrices
$H_m^{1/2}P_jH_m^{1/2}$ and $H_m^{1/2}Q_jH_m^{1/2}$ are their
exact Euclidean isometric conjugates.

Differentiating the norm of the monic polynomial gives
$\dot\omega_j=p_j^*\dot H_jp_j$: its derivative has degree less
than $j$, and both cross terms vanish by its original orthogonality.
Differentiating $J_jR_j=I$ makes $\dot R_j$ a boundary, so the
same expansion gives $\dot G_j=R_j^*\dot H_jR_j$.
The determinant derivative follows by expanding the determinant
multilinearly in its columns, or by differentiating
$\det(G_j+t\dot G_j)=\det G_j\det(I+tG_j^{-1}\dot G_j)$.
Thus
\[
 \frac d{dt}\log\omega_j=\Tr(P_jE),\qquad
 \frac d{dt}\log V_j=\Tr(Q_jE).                             \tag{EW.31}
\]
All these formulas use the original coordinates before the final
cyclic trace.

Put $\mathscr L=\log\mathscr R\ge0$ and $s=\min(q,r)$.
The endpoint derivatives are exactly
\[
 \begin{split}
 (\log U)'&=\Tr((P_m-P_n)E),\\
 \mathscr L'&=\Tr((Q_{n-1}+Q_n-Q_{m-1}-Q_m)E),\\
 \mathcal E_{n,r}'&=\frac1{2r}\left\{
 \mathcal E_{n,r}(\log U)'
 +U^{1/(2r)}\cosh(\mathscr L/(2r))\mathscr L'\right\}.
 \end{split}                                                   \tag{EW.32}
\]
These equalities also hold at $\mathscr R=1$. They identify the
exact paired sensitivity rather than separate errors of large source
and relation determinants.

\begin{theorem}[Intrinsic finite sensitivity]
Throughout every such positive continuously differentiable family,
\[
 \boxed{|(\log U)'|\le\mathfrak o_H(E),\qquad
 |\mathscr L'|\le\mathfrak o_H(E)
             \min\{2s,\sqrt{2s\mathscr L}\}.}                \tag{EW.33}
\]
For any two positive source matrices $H^{(0)},H^{(1)}$ in the
same original coordinates, let
\[
 \kappa=\frac{\lambda_{\max}((H^{(0)})^{-1/2}H^{(1)}(H^{(0)})^{-1/2})}
                   {\lambda_{\min}((H^{(0)})^{-1/2}H^{(1)}(H^{(0)})^{-1/2})}
 \ge1.
                                                               \tag{EW.34}
\]
Their original window values satisfy
\[
 \boxed{|\log(U^{(1)}/U^{(0)})|\le\log\kappa,\qquad
 |\sqrt{\mathscr L^{(1)}}-\sqrt{\mathscr L^{(0)}}|
                  \le\sqrt{s/2}\log\kappa.}                 \tag{EW.35}
\]
In particular these sensitivity bounds contain no independent
source-size determinant exponent. They vanish for any scalar change
$H^{(1)}=tH^{(0)}$, $t>0$, exactly as the original endpoint does.
\end{theorem}
\begin{proof}
Let $\Lambda=H_m^{-1/2}EH_m^{-1/2}$. Subtracting
$\tfrac12(\lambda_{\max}\Lambda+\lambda_{\min}\Lambda)I$
leaves operator norm $\mathfrak o_H(E)/2$. This scalar subtraction
does not change either trace in (EW.32): the two monic-line
projections each have trace one and the four least-lift projections
each have trace $q$, with the displayed signs summing to zero.

The two monic orthogonal lines at degrees $n<m$ are orthogonal
in the common ambient source. Their projection difference has
eigenvalues $1,-1$ and zeros, hence trace norm two. For any
self-adjoint $A$ and $B$, diagonalizing $A$ gives
$|\Tr(AB)|\le\sum_i|\lambda_i(A)|\,\|B\|$.
This proves the first bound in (EW.33).

For the second trace split its four projections as
$(\Pi_{n-1}-\Pi_{m-1})+(\Pi_n-\Pi_m)$. Both index gaps
are exactly $r$. Put
$\ell_1=\log(V_{n-1}/V_{m-1})$ and
$\ell_2=\log(V_n/V_m)$, so $\ell_1+\ell_2=\mathscr L$.
Equation (EW.28), the triangle inequality for the trace norm, and
$\sqrt{\ell_1}+\sqrt{\ell_2}\le\sqrt{2\mathscr L}$ give
\[
 \|\Pi_{n-1}+\Pi_n-\Pi_{m-1}-\Pi_m\|_1
 \le 2\sqrt{2s\mathscr L},\qquad
 \|\Pi_{n-1}+\Pi_n-\Pi_{m-1}-\Pi_m\|_1\le4s.
\]
Multiplying either bound by $\mathfrak o_H(E)/2$ proves the
second bound of (EW.33).

For a continuously differentiable path and $\varepsilon>0$, the exact bound
$|\mathscr L'|\le\mathfrak o_H(E)\sqrt{2s\mathscr L}$ implies
\[
 \left|\frac d{dt}\sqrt{\mathscr L+\varepsilon}\right|
 \le\sqrt{s/2}\,\mathfrak o_H(E).
\]
Integrate on the compact path and let $\varepsilon$ decrease to
zero. This proves the square-root Lipschitz assertion along a path,
including points at which $\mathscr L=0$; no division by zero is
used.

Use the straight positive segment
$H(t)=(1-t)H^{(0)}+tH^{(1)}$, $0\le t\le1$.
Let $\ell_i>0$ be the eigenvalues of
$(H^{(0)})^{-1/2}H^{(1)}(H^{(0)})^{-1/2}$.
Congruence by $(H^{(0)})^{1/2}$ and then its displayed orthonormal
eigenbasis gives the generalized eigenvalues of $(E,H(t))$ exactly
as
\[
 \frac{\ell_i-1}{1+t(\ell_i-1)}.
\]
These are the eigenvalues of $H(t)^{-1/2}EH(t)^{-1/2}$:
both follow from the same equation $Ex=\lambda H(t)x$.
The displayed function is increasing in $\ell_i$, with derivative
$[1+t(\ell_i-1)]^{-2}$. Hence
\[
 \int_0^1\mathfrak o_{H(t)}(E)\,dt
 =\log\ell_{\max}-\log\ell_{\min}=\log\kappa.
\]
Integrating the already proved differential inequalities proves
(EW.35). Finally direct substitution in (EW.11) gives
$R_N(tH)=R_N(H)$, $G_N(tH)=tG_N(H)$, and
$\omega_j(tH)=t\omega_j(H)$. Thus $U$ and $\mathscr R$ are
exactly unchanged by that explicitly stated common scaling. The
source mass remains $t\mu_h^k$ in the scaled comparison; it was
never reassigned the value one.
\end{proof}

The independent projection review supplied a stronger exact version
from the same angles. For $s\ge1$, it is
\[
 \boxed{|\mathscr L'|\le
 2s\,\mathfrak o_H(E)\sqrt{1-e^{-\mathscr L/(2s)}}.}         \tag{EW.35a}
\]
To prove it, list $x_i=-\log\sigma_i^2$ for each of the two
projection pairs used above, omitting their zero terms and then
appending zeros to make exactly $2s$ entries. Their sum is
$\mathscr L$. The exact trace-norm formula and the scalar
cancellation prove
$|\mathscr L'|\le\mathfrak o_H(E)\sum_i f(x_i)$, where
$f(x)=\sqrt{1-e^{-x}}$. For $x>0$,
\[
 f''(x)=-\frac{e^{-x}}{2\sqrt{1-e^{-x}}}
        -\frac{e^{-2x}}{4(1-e^{-x})^{3/2}}<0.
\]
Continuity at zero extends concavity there: apply the concavity
inequality to positive arguments increased by $\varepsilon$ and
then let $\varepsilon\downarrow0$. Finite Jensen on the full
$2s$-term list proves (EW.35a). Since $1-e^{-x}\le x$ and
$1-e^{-x}\le1$, it implies both bounds on $|\mathscr L'|$
in (EW.33).

Define the explicit increasing map on the actual loss variable by
\[
 \Phi_s:[0,\infty)\longrightarrow[0,\infty),\quad
 \Phi_s(L)=2\operatorname{arcosh}(e^{L/(4s)}),\qquad
 \Phi_s^{-1}(v)=4s\log\cosh(v/2).
\]
Here $\operatorname{arcosh}x=\log(x+\sqrt{x^2-1})$ for
$x\ge1$, so both maps and their inverses have specified real
branches. For $L>0$,
$\Phi_s'(L)=[2s\sqrt{1-e^{-L/(2s)}}]^{-1}$.
Apply (EW.35a) to $\Phi_s(\mathscr L+\varepsilon)$.
Its derivative in absolute value is at most
$\mathfrak o_H(E)\sqrt{(1-e^{-\mathscr L/(2s)})/
(1-e^{-(\mathscr L+\varepsilon)/(2s)})}\le\mathfrak o_H(E)$.
Integrating and letting $\varepsilon\downarrow0$, then using
the exact straight-segment length already proved, gives
\[
 \boxed{|\Phi_s(\mathscr L^{(1)})-
               \Phi_s(\mathscr L^{(0)})|\le\log\kappa.}      \tag{EW.35b}
\]
This conclusion includes zero losses by continuity. For $s=0$,
the only loss is zero and define the one-point map $\Phi_0(0)=0$;
the projection and endpoint comparisons are their zero identities.

For $s\ge1$, the stronger enclosure itself has an exact algebraic form.
Given any $\rho\ge\kappa$ and the original $\mathscr R^{(0)}\ge1$,
put
\[
 \begin{gathered}
 z_0=(\mathscr R^{(0)})^{1/(4s)}
       +\sqrt{(\mathscr R^{(0)})^{1/(2s)}-1},\qquad
 z_-=\max\{1,z_0/\sqrt\rho\},\quad z_+=\sqrt\rho\,z_0,\\
 \boxed{\left(\frac{z_-+z_-^{-1}}2\right)^{4s}
       \le\mathscr R^{(1)}\le
       \left(\frac{z_++z_+^{-1}}2\right)^{4s}.}
 \end{gathered}                                               \tag{EW.35c}
\]
Indeed $z_0=\exp(\Phi_s(\mathscr L^{(0)})/2)$,
(EW.35b) encloses its counterpart between $z_-$ and $z_+$,
and $z\mapsto(z+z^{-1})/2$ is increasing for $z\ge1$.
The identity $\mathscr R=\cosh(\Phi_s(\log\mathscr R)/2)^{4s}$
proves (EW.35c). At the same time
$U^{(0)}/\rho\le U^{(1)}\le\rho U^{(0)}$ by (EW.35).
Thus monotonicity in (EW.20) supplies an endpoint enclosure using
only nonnegative real roots, rational operations and integer powers when
the central $U,\mathscr R$ and a bound $\rho$ are rational.
For example any proved rational Loewner comparison
$aH^{(0)}\preceq H^{(1)}\preceq bH^{(0)}$, $0<a\le b$,
gives $\rho=b/a$: the extreme Rayleigh quotients lie between
$a$ and $b$. This is the exact map from a positive-matrix
certificate to the sharper window certificate, with the original
mass and relation unchanged.

The square-root assertion is a Lipschitz assertion, and does not
assert ordinary differentiability at every zero of $\mathscr L$.
For an exact auxiliary illustration take $q=1$, $\chi=S$, $n=1$,
$r=1$, and the positive original-coordinate matrix
\[
 H_2(t)=\begin{pmatrix}1&0&t\\0&1&0\\t&0&1\end{pmatrix},
 \qquad |t|<1.
\]
Its leading pivots are $1,1,1-t^2$, its remainder is evaluation
at $S=0$, and (EW.11) gives $V_0=V_1=1$, $V_2=1-t^2$.
Thus $\mathscr L=-\log(1-t^2)$ and
$\sqrt{\mathscr L}/|t|\to1$ as $t\to0$. Its square root
has a corner. This positive-form example belongs exactly to the
finite family domain of (EW.33)--(EW.35); it is not claimed to be
an arithmetic moment source. The regularized proof above includes it.

For completeness, (EW.35) has a directly usable endpoint enclosure.
Writing $d_H=\log\kappa$ and
$b=\sqrt{s/2}\,d_H$, it proves
\[
 \begin{split}
 e^{-d_H}U^{(0)}&\le U^{(1)}\le e^{d_H}U^{(0)},\\
 \bigl(\max\{0,\sqrt{\mathscr L^{(0)}}-b\}\bigr)^2
 &\le\mathscr L^{(1)}
 \le(\sqrt{\mathscr L^{(0)}}+b)^2.
 \end{split}                                                   \tag{EW.36}
\]
The function $U^{1/(2r)}\sinh(\mathscr L/(2r))$ is increasing
in each positive $U$ and nonnegative $\mathscr L$, as direct
differentiation shows. Substitution of the two endpoints in (EW.36)
therefore bounds the actual endpoint expression. This is a finite
stability theorem, including its exact $s=\min(q,r)$ dependence.

\subsection{Rational interval certificates for the original input}

Here is an explicit finite algorithm which uses enclosures of the
original coefficients; it does not replace them by an arbitrary moment
matrix. For clarity first suppose the constants in (EW.5) and the
polynomial tests in (EW.10) are exact. Given real intervals
$|c_{h,j}-\widehat c_j|\le e_j$ for $0\le j\le2m$, set
\[
 \begin{split}
 \widehat A(z)&=\sum_{j=0}^{2m}(\alpha)_j\widehat c_jz^j,\\
 A_+(z)&=\sum_{j=0}^{2m}(\alpha)_j|\widehat c_j|z^j,
 \qquad E_+(z)=\sum_{j=0}^{2m}(\alpha)_je_jz^j,\\
 \widehat d_l&=[z^l]\widehat A(z)^k,\quad
 e_l^{(k)}=[z^l]\{(A_+(z)+E_+(z))^k-A_+(z)^k\}.
 \end{split}                                                   \tag{EW.37}
\]
The multinomial expansion groups each difference of products by
the nonempty set of factors carrying an error. Taking absolute values
term by term proves $|d_l-\widehat d_l|\le e_l^{(k)}$.
Thus the Hermitian central matrix and entry radii
\[
 \widehat H_{ab}=C_k\sum_l l!L_{ab,l}\widehat d_l,\qquad
 \varepsilon_{ab}=C_k\sum_l l!|L_{ab,l}|e_l^{(k)}            \tag{EW.38}
\]
enclose the actual source. Rational upper bounds for the moduli give
rational radii. An uncertain real constant, including $C_k$, is
propagated by interval multiplication before the final matrix is
reported. For $|x-\widehat x|\le e_x$ and
$|y-\widehat y|\le e_y$, the exact elementary rule is
$|xy-\widehat x\widehat y|\le|\widehat x|e_y+
|\widehat y|e_x+e_xe_y$. Repeated use handles every finite
product in (EW.8)--(EW.10).

The original relation need not have rational coefficients. To retain
it exactly while producing rational certificates, use its quotient-first
coordinate map at each endpoint degree $N$:
\[
 A_N=[E_N,B_N]:\C^q\oplus\C^{N-q+1}\longrightarrow\C^{N+1},
 \quad\det A_N=1,\qquad\widetilde H_N=A_N^*H_NA_N.
                                                               \tag{EW.39}
\]
The first component of $A_N^{-1}P$ is exactly $J_NP$ by monic
division. Supply rational enclosures of the real and imaginary parts
of the actual $\chi$ coefficients when they are not exact rational
numbers. Ordinary complex interval addition, multiplication and
conjugation in (EW.39) give rational entry enclosures of the actual
$\widetilde H_N$. For a complex interval a modulus radius is
bounded rationally, for example by the sum of its two coordinate
radii. Conjugation changes the sign of the imaginary coordinate;
it is applied to the actual first factor. This procedure never
changes $\chi$, its degree, its monicity or the original quotient
coordinates. Intervals for $\lambda$ and its positive constants,
if required, are propagated in the same finite operations and in the
specified recurrence (EW.6).

Here is a rational matrix certificate valid for either $H_m$ or
any of these $\widetilde H_N$. Choose a Hermitian rational centre
$\widehat M$ and symmetric rational radii $\varepsilon_{ab}$
enclosing its actual error. Choose any positive rational $v_a$ and
form
\[
 \Delta=\diag\left(\sum_b\varepsilon_{ab}\frac{v_b}{v_a}\right),
 \qquad M^-=\widehat M-\Delta,\quad M^+=\widehat M+\Delta.
                                                               \tag{EW.40}
\]
The pairwise inequality
$2|x_a||x_b|\le(v_b/v_a)|x_a|^2+(v_a/v_b)|x_b|^2$
proves $M^-\preceq M\preceq M^+$. Rational complex arithmetic
now eliminates each leading pivot by the congruence
\[
 \begin{pmatrix}a&b^*\\b&C\end{pmatrix}
 \longmapsto\diag(a,C-bb^*/a),\qquad
 (x,y)\longmapsto(x-a^{-1}b^*y,y).                           \tag{EW.41}
\]
Recording strictly positive rational pivots proves $M^->0$ by
induction and hence also $M^+>0$. A nonpositive pivot rejects that
particular proposed certificate. Every actual matrix is already
positive by its source integral and, for (EW.39), its invertible
congruence. At a fixed finite degree, convergent valid entry intervals
eventually give positive lower matrices: the least unit-sphere value
of the actual quadratic form is positive by compactness, and the
finite sum of entry errors tends to zero. This proves the finite
certificate's convergence without assuming or claiming that a
growing-degree arithmetic input has already been enclosed.

For the two norm endpoints take the corresponding principal
submatrices of the certified $H_m^-,H_m^+$. Define
$\omega_j^-,\omega_j^+$ by their exact monic Schur complements.
The squared norm of any fixed monic polynomial is between its two
enclosing forms. Minimizing over the same affine monic set proves
\[
 0<\omega_j^-\le\omega_j\le\omega_j^+,
 \qquad j=n,m.                                               \tag{EW.42}
\]
For each of $N=n-1,n,m-1,m$, use (EW.39)--(EW.41) and split
its first $q$ and last $N-q+1$ coordinates. The Schur complements
of the latter block in $\widetilde H_N^-,\widetilde H_N^+$
give rational positive forms $G_N^-,G_N^+$. Their affine fibre
has fixed first coordinate $x$, by the exact map (EW.39).
Minimization therefore proves
\[
 G_N^-\preceq G_N\preceq G_N^+,\qquad
 0<V_N^-:=\det G_N^-\le V_N\le\det G_N^+=:V_N^+.
                                                               \tag{EW.43}
\]
For the determinant inequality, conjugate by the square root of
the smaller positive form and multiply the resulting eigenvalues,
each at least one. At $N=q-1$ the boundary block is empty and
the Schur complement is the whole $q$-dimensional form.

All numbers in the following enclosure are rational and positive:
\[
 \begin{aligned}
 U_-&=\omega_m^-/\omega_n^+,&
 U_+&=\omega_m^+/\omega_n^-,\\
 \mathscr R_-&=\max\left\{1,
 \frac{V_{n-1}^-V_n^-}{V_{m-1}^+V_m^+}\right\},&
 \mathscr R_+&=\frac{V_{n-1}^+V_n^+}{V_{m-1}^-V_m^-}.
 \end{aligned}                                                 \tag{EW.44}
\]
They satisfy $U_-\le U\le U_+$ and
$1\le\mathscr R_-\le\mathscr R\le\mathscr R_+$.
The use of the lower bound one follows from the proved nested
minimization, rather than an interval rounding choice.

Positive roots require no floating logarithm certificate. For any
positive rational $x$ and integer $2r$, find rational numbers
$l,u>0$ such that $l^{2r}\le x\le u^{2r}$. Doubling upper
or lower powers of two supplies an initial interval, and ordinary
rational bisection makes it arbitrarily narrow. Each decision is a
positive-denominator integer comparison. Strict monotonicity of
$t\mapsto t^{2r}$ for $t>0$ proves the enclosure and convergence.
Apply this algorithm to
\[
 U_-\mathscr R_-,\quad U_-/\mathscr R_-,\quad
 U_+\mathscr R_+,\quad U_+/\mathscr R_+.
\]
Write their certified root intervals as $[l_1,u_1],\ldots,[l_4,u_4]$.
The function in (EW.20) is increasing in $U>0$ and
$\mathscr R\ge1$: its $U$ derivative is its nonnegative value
times $1/(2rU)$, and its $\mathscr R$ derivative is
$U^{1/(2r)}\cosh(\log\mathscr R/(2r))/(2r\mathscr R)>0$.
Therefore the complete rational endpoint certificate is
\[
 \boxed{\max\{0,(l_1-u_2)/2\}
             \le\mathcal E_{n,r}\le(u_3-l_4)/2.}             \tag{EW.45}
\]
Alternatively (EW.22) gives a positive numerator and denominator
certificate from the same root brackets. Every matrix pivot, integer
power comparison, relation coefficient and source mass used in this
algorithm is inspectable. It makes no claim for coefficient intervals
which have not been provided.

\subsection{Boundary cases and the retained cochain map}

At $r=1$, the block has two columns, $W_r=w$, and
$\mathscr R=\mathscr R_2$ because the two copies of
$\mathscr R_1$ cancel. Equations (EW.18)--(EW.22) reduce exactly
to the two-degree formula GMT.24--27. At $n=q$, the base source
degree is $q-1$ and its boundary dimension is zero; its quotient
metric is simply the source Gram. All formulas above still use
positive matrices of their specified sizes. No degree $q-2$
full-quotient form is inserted.

The actual positive polynomial source has an additional strictness
property for $q\ge1$: every two-degree quotient-volume ratio is
strictly greater than one. Here is the complete algebraic reason.
Let $\pi_j(u)$ be the real monic orthogonal family of the original
positive real-$u$ measure. Multiplication by $u$ is self-adjoint.
Pairing $u\pi_j$ with lower orthogonal polynomials and using their
degrees gives
\[
 u\pi_j=\pi_{j+1}+b_j\pi_j+a_j\pi_{j-1},\qquad
 b_j\in\R,\quad a_j=\omega_j/\omega_{j-1}>0\ (j\ge1).
\]
For $j=0$ the last term is absent. The exact monic phase map
$p_j(S)=i^j\pi_j((S-c)/i)$ has the same squared norm, because
$|i^j|=1$, and gives
\[
 S p_j=p_{j+1}+(c+ib_j)p_j-a_jp_{j-1}. 
\]
If a nonconstant $\chi$ divided $p_j$ and $p_{j+1}$, this
recurrence and $a_j\ne0$ would make it divide $p_{j-1}$.
Induction would make it divide $p_0=1$, a contradiction.
For the two appended columns at degrees $N,N+1$ over
$\mathcal P_{N-1}$, the residuals are $p_N$ and
$p_{N+1}+z p_N$, where the original coefficient is exactly
$z=\langle p_N,S^{N+1}\rangle_H/\omega_N$.
Indeed their difference from those polynomials has smaller degree,
and both residuals are orthogonal to $\mathcal P_{N-1}$.
If their full quotient residual map $F_2$ were zero, both $p_N$
and $p_{N+1}$ would have zero remainder, contradicting the
divisibility argument. Thus $Z_2=F_2^*G_{N-1}F_2$ is a
nonzero positive semidefinite matrix, and (EW.17) gives
$V_{N-1}/V_{N+1}>1$ for every $N\ge q$.
The finite product
\[
 \mathscr R=\prod_{j=0}^{r-1}
                  \frac{V_{n+j-1}}{V_{n+j+1}}>1
 \qquad(q\ge1,n\ge q,r\ge1)                                \tag{EW.46a}
\]
is its exact overlapping telescope. This proves strict positivity
of the actual $\mathcal E_{n,r}$ and of $\mathscr L$ at each
admitted nonempty arithmetic window; it supplies no uniform lower
bound as degrees or packets grow. The more general positive-form
sensitivity theorem includes the zero-loss boundary separately.

The empty packet is a separate declared case: $h=1$, $\chi=1$,
$q=0$, and $\mathcal C=0$. The analytic density and mass remain
those of (EW.1). Put $V_N=1$ even at the empty degree $N=-1$;
all quotient matrices, projections and $F_j$ have their empty
dimensions. For $n\ge0,r\ge1$ the source norms exist,
$\mathscr R=1$, and $\mathcal E_{n,r}=0$. If $n=0$, use the
empty base source $H_{-1}$, with determinant one and no inverse
operation; the block $W=H_r$ then has its original monomial
coordinates. Its first entry is the full mass. The projection and
volume sensitivity statements become their zero-map identities,
while the source-norm sensitivity remains the proved rank-two
calculation.

Finally the projection comparison (EW.26) has an explicit original
theta representative. For $a<b$ and $x\in\mathcal C$, monic
division gives a unique polynomial $\kappa_x\in\mathcal P_{b-q}$
with
\[
 \iota_aR_ax-\iota_bR_bx=\chi\kappa_x.                       \tag{EW.46}
\]
Its coefficients are obtained by the same exact division, so the
map $x\mapsto\kappa_x$ is linear. Let
$\mathcal V_{h,k}P=P(D_1+\cdots+D_k)F_h^{\otimes k}$, where
$D=-x\partial_x$, $\mathcal MF_h=v_h$, and
$h(D)F_h=\Theta\phi_0$. Successive division in
$s_1,\ldots,s_k$ supplies specific $Q_i$ with
$\chi(\sum s_i)=\sum_i h(s_i)Q_i(\mathbf s)$.
The final tensor remainder vanishes because its monomials form a
basis of $(\C[s]/(h))^{\otimes k}$ and $\chi$ annihilates
the original sum. Thus the exact cochain of degree $k-1$
\[
 \mathcal P_\chi(\kappa_x)=
 \sum_{i=1}^k(-1)^{i-1}Q_i(\mathbf D)\kappa_x(\textstyle\sum D_j)
 \bigl(F_h^{\otimes(i-1)}\otimes\phi_0\otimes
                         F_h^{\otimes(k-i)}\bigr)            \tag{EW.47}
\]
satisfies
\[
 d_{\rm tot}\mathcal P_\chi(\kappa_x)
   =\mathcal V_{h,k}(\iota_aR_ax-\iota_bR_bx).                \tag{EW.48}
\]
Indeed the tensor differential at position $i$ contributes
$(-1)^{i-1}$ because its preceding factors have cochain degree
one; this cancels the written sign, and $h(D_i)F_h=\Theta\phi_0$
then proves the equality term by term. The arithmetic jet of both
representatives is the same $\eta x$ of (EW.3). In the original
support carrier $\{\tau\}\sqcup\bigsqcup_\ell(\{\ell\}\times E_\ell)$,
a fibrewise linear map sends $(\ell,v)$ to $(\ell,L_\ell v)$
and $\tau$ to $\tau$. Thus the displayed boundary maps to
$e_\ell=(\ell,0)$ while external absence remains $\tau$.
The endpoint's representative changes, projection angles, and volume
changes have therefore all been related to the original typed maps
and their full theta primitives.

\subsection{Verification scope and exact auxiliary calibration}

The companion checker uses the explicitly auxiliary weight
\[
 \begin{gathered}
 \lambda=1,\quad k=2,\quad c=1,\quad c_1=1/2,\quad C_2=1/4,\\
 B(t)=1+t^2/4,\quad w=B r_1,\quad\chi(S)=S^2-2S-1.
 \end{gathered}                                               \tag{EW.49}
\]
This is a specified positive $L^2(r_1)$ multiplier and the same
finite coefficient-to-Gram domain as (EW.8)--(EW.10). It is not
asserted to be an actual zeta packet. Its exact connection to the
arithmetic problem is the complete sequence of maps (EW.10),
(EW.11), (EW.14)--(EW.23), and (EW.39)--(EW.45), with its
auxiliary inputs written explicitly.

For this fixture $b_2^{(1)}=t^2-2$, so
$c_{B,0}=3/2$, $c_{B,2}=1/4$, and every other coefficient is
zero. Consequently $\sum_j(2)_jc_{B,j}z^j=3/2+3z^2/2$ and
its square has coefficients $d_0=9/4,d_2=9/2,d_4=9/4$.
Its one-factor mass is $3/4$, the literal sum mass is $9/16$,
and the reference sum mass is $1/4$. The checker independently
constructs Gamma moments from the monic recurrence, inserts
$1+t^2/4$ in each factor, and expands every power of
$t_1+t_2$ with its full binomial coefficient. It also constructs
the coefficient-side moments from (EW.9); agreement compares
two exact finite calculations, including the tensor cross terms.

In this auxiliary example gamma recurrence gives a further exact
all-degree check. Since
$|\Gamma(2+it/2)|^2=(1+t^2/4)|\Gamma(1+it/2)|^2$, one has
$Br_1=r_2$. The displayed masses give
$c_2=3/4$, $c_4=315/8$, and hence
\[
 m=r_2*r_2=r_4/70,\qquad
 r_1*r_1=r_2/3,\qquad
 \omega_j=\frac9{16}j!(8)_j,\qquad
 U=(n+1)_r(n+8)_r.                                         \tag{EW.50}
\]
The first two equalities follow from the literal convolution law
(EW.7), with coefficients $(3/4)^2/(315/8)=1/70$ and
$(1/2)^2/(3/4)=1/3$. The monic original-coordinate map is
$p_j(S)=i^jb_j^{(4)}((S-1)/i)$, whose squared norm is exactly
the third quantity by (EW.7). Taking the stated endpoint ratio
proves the final formula.

The exact selected windows computed in the accompanying rational
receipt are
\[
\begin{array}{c|c|c}
 (n,r)&U&\mathscr R\\\hline
 (2,1)&30&20191/4860\\
 (2,2)&1320&3073295611/216513000\\
 (3,2)&2640&295913127219469/36104825790000
\end{array}                                                   \tag{EW.51}
\]
Here is a full independent finite derivation of the entries in that
table, using the original quotient kernel rather than a Schur block.
Let $x=S-1$ solely in the stated quotient-coordinate map
\[
 A:\C^2_{(1,x)}\longrightarrow\C^2_{(1,S)},\qquad
 (a,b)\longmapsto(a-b,b),\qquad
 A=\begin{pmatrix}1&-1\\0&1\end{pmatrix},\quad\det A=1.
\]
It takes the identical relation to $x^2-2$, so its residue equation
is exactly $x^2=2$. It preserves the determinant of the quotient
form by the congruence $G_N^x=A^*G_NA$. The monic source family
in (EW.50) obeys
$p_{j+1}=xp_j+j(j+7)p_{j-1}$, by applying the exact phase map to
(EW.6). Orthogonal expansion of the original source inverse in
(EW.11) proves the finite quotient kernel formula
\[
 (G_N^x)^{-1}=\sum_{j=0}^N
     \frac{[p_j]_{(1,x)}[p_j]_{(1,x)}^*}{\omega_j}.
                                                               \tag{EW.52}
\]
Indeed the matrix of the monic orthogonal columns in the original
source is invertible triangular with diagonal one, takes $H_N$
to $\diag(\omega_0,\ldots,\omega_N)$, and sends the reduction
columns to their written remainders. Substituting that inverse
congruence into $J_NH_N^{-1}J_N^*$ gives (EW.52) term by term.
The recurrence and $x^2=2$ give the complete required list
\[
\begin{array}{c|c|c}
 j &[p_j]_{(1,x)}&\omega_j\\\hline
 0&(1,0)&9/16\\
 1&(0,1)&9/2\\
 2&(10,0)&81\\
 3&(0,28)&2430\\
 4&(356,0)&106920\\
 5&(0,1588)&6415200
\end{array}                                                   \tag{EW.53}
\]
For example the successive remainders are obtained from
$x^2+8=10$, $10x+18x=28x$, $28x^2+300=356$, and
$356x+44\cdot28x=1588x$. The positive kernels from (EW.52)
are diagonal in this explicitly transported quotient basis, with
\[
\begin{array}{c|c|c|c}
 N&((G_N^x)^{-1})_{00}&((G_N^x)^{-1})_{11}&V_N\\\hline
 1&16/9&2/9&81/32\\
 2&244/81&2/9&729/488\\
 3&244/81&662/1215&98415/161528\\
 4&56102/13365&662/1215&16238475/37139524\\
 5&56102/13365&376069/400950&2679348375/10549111519
\end{array}                                                   \tag{EW.54}
\]
Each diagonal entry is the indicated finite sum of rational terms
in (EW.52), and $V_N$ is the reciprocal of their product because
$\det A=1$. Substituting (EW.53)--(EW.54) in (EW.18) gives all
three rows of (EW.51) by positive-denominator multiplication.

There is also an executed coefficient-box certificate, before any
matrix error is chosen. Put $\varepsilon=10^{-12}$ and retain
$n=2,r=1,m=3$, the same reference mass and the same $\chi$.
Allow the complete specified one-factor input box
\[
 |c_{B,0}-3/2|\le\varepsilon,\qquad
 |c_{B,2}-1/4|\le\varepsilon,\qquad
 c_{B,j}=0\quad(j=1,3,4,5,6).                               \tag{EW.55}
\]
Every point of this box is realized by the specified positive multiplier
$B_{a,b}(t)=a+bt^2$, where $a=c_{B,0}-2c_{B,2}$ and
$b=c_{B,2}$. They obey $a\ge1-3\varepsilon>0$ and
$b\ge1/4-\varepsilon>0$. Its coefficient formula follows by
$a+bt^2=(a+2b)b_0^{(1)}+b b_2^{(1)}$.
The interior input $a=1+\varepsilon$, $b=1/4-\varepsilon/4$
has one-factor mass $3/4+\varepsilon/4$ and sum mass
$(3/4+\varepsilon/4)^2$. Thus neither the box nor its test input
replaces the zeroth coefficient by a unit mass.

For (EW.55), the complete coefficient-power error formula (EW.37)
is exactly
\[
 \begin{gathered}
 e_0^{(2)}=3\varepsilon+\varepsilon^2,\quad
 e_2^{(2)}=21\varepsilon+12\varepsilon^2,\quad
 e_4^{(2)}=18\varepsilon+36\varepsilon^2,\\
 e_1^{(2)}=e_3^{(2)}=e_5^{(2)}=e_6^{(2)}=0.
 \end{gathered}                                               \tag{EW.56}
\]
Indeed the central polynomial is $3/2+(3/2)z^2$ and its error
polynomial is $\varepsilon+6\varepsilon z^2$; expanding the
full square difference gives the three displayed coefficients.
The upper corner of the coefficient box attains all three positive
differences, so its quadratic error terms cannot be omitted.

Expanding all sixteen original tests $(1-iu)^a(1+iu)^b$,
$0\le a,b\le3$, in the monic Gamma basis of (EW.6), then using
(EW.38)--(EW.40) with four weights equal to one, gives
\[
 \widehat H_3=\frac1{16}\begin{pmatrix}
 9&9&-63&-207\\9&81&81&-1863\\
 -63&81&2025&2025\\-207&-1863&2025&93393
 \end{pmatrix},
\]
\[
 \begin{aligned}
 \Delta_{00}&=54\varepsilon+28\varepsilon^2,&
 \Delta_{11}&=(1029\varepsilon+873\varepsilon^2)/2,\\
 \Delta_{22}&=(2109\varepsilon+1773\varepsilon^2)/2,&
 \Delta_{33}&=29754\varepsilon+32878\varepsilon^2.
 \end{aligned}                                                 \tag{EW.57}
\]
These diagonal sums retain the absolute value of every test
coefficient. The real coefficients occur at even degrees and the
imaginary coefficients at odd degrees; the latter pair exactly to
zero against every member of the specified even coefficient box.
The complete complex expansions, including those zero pairings,
are recorded in the certificate.

Here is an elementary exact positive-definiteness certificate for
the whole box. Rational multiplication verifies
\[
 \widehat H_3^{-1}=\frac1{2430}\begin{pmatrix}
 8019&-1863&351&-27\\-1863&1501&-147&29\\
 351&-147&39&-3\\-27&29&-3&1
 \end{pmatrix},\qquad \Tr\widehat H_3^{-1}=956/243<4.
\]
The central form is positive by its actual measure, or by the
four original monic norms in (EW.53). Its inverse is positive,
so each of its eigenvalues is at most its trace. Consequently
$\widehat H_3\succ I_4/4$.
Substitution of $\varepsilon=10^{-12}$ into the four full
polynomials (EW.57) gives $0<\Delta_{jj}<10^{-6}$ by integer
comparison. Hence
\[
 \widehat H_3-\Delta\succ(1/4-10^{-6})I_4\succ0,
 \qquad
 \widehat H_3-\Delta\preceq H_3\preceq\widehat H_3+\Delta
                                                               \tag{EW.58}
\]
for every actual multiplier in (EW.55). This proves the positive
matrix enclosure before its inverses are evaluated. The accompanying
receipt also records each exact rational elimination pivot separately.

Use the original $J_N$ and $\chi$ at $N=1,2,3$ in
(EW.42)--(EW.45). The resulting rational interval for every
endpoint value in this whole coefficient box is
\[
 \boxed{\frac{10247872526420120792621067}{2^{81}}
 \le\mathcal E_{2,1}\le
 \frac{10247872551311149445602506}{2^{81}}.}                  \tag{EW.59}
\]
The four endpoint arguments in (EW.45) are computed from the
written rational matrices $\widehat H_3\pm\Delta$ by their
monic and quotient Schur formulas; the receipt gives every one
of these exact rational values and its two squared-root comparisons.
Thus (EW.59) is certified by integer powers and the minimum
principle, with no rounding assertion about a transcendental
logarithm. The interior positive multiplier and all four box
corners are independently reconstructed by the literal one-factor
moment recurrence and tensor binomial expansion, and pass the
same original-coordinate entry bounds.

All original-$S$ Gram entries are computed from
$(1-iu)^a(1+iu)^b$. The checker's window certificates include
the exact reduction matrix, source norms, quotient metrics,
the $r+1$ Schur block and all its prefixes, both cancellations
in (EW.21), positive rational matrix envelopes, and the integer
power comparisons for (EW.45). The written proof establishes
the all-degree statements; each executed finite check is reported
separately in the accompanying receipt. The base fixture executes
416 distinct mathematical checks across fifteen windows; its normal
and optimized runs pass and six changed-formula runs in each mode
are rejected after executing the complete suite. The separate
supplement executes 329 mathematical checks and two dependency-hash
checks; its four changed-formula runs in each mode are also rejected
after the complete suite, including the omitted quadratic coefficient
error at an attained box corner. The independent remainder-kernel
calibration executes the six scalar comparisons in (EW.51).
These counts describe their separate scopes; repeated modes and
changed-formula runs do not add new mathematical checks. Numerical
examples do not stand in for the arithmetic input or a growing-window bound.


\endgroup
